% =====================================================================
%  Section: Related Work
% =====================================================================
\section{Related Work}
\label{sec:related}

\subsection{Image Segmentation for Artistic Content}

Instance segmentation in natural images has matured rapidly with the YOLO
family~\cite{jocher2023yolov8,wang2024yolov9,wang2024yolov10,tian2025yolov12}
and the Segment Anything Model
(SAM)~\cite{kirillov2023sam,ravi2024sam2}. However, paintings pose unique
challenges: brushstroke textures confound edge detectors, semantic boundaries
blend intentionally in impressionist works, and foreground-background separation
lacks the depth cues present in photographs. Specialized art-domain
adaptations have begun to emerge---PLDS-YOLO~\cite{chen2025pldsyolo} adapts
YOLOv8 for mural paint-loss detection, and art-specific segmentation
datasets~\cite{saleh2016metric} provide annotated painting regions---but a
general-purpose painting segmentation pipeline suitable as the front-end of
a 3D animation framework remains absent. Canvas3D bridges this gap by
fine-tuning a YOLO-family detector on a curated art-segmentation dataset,
enabling reliable foreground-object extraction from diverse painting styles.

\subsection{3D Gaussian Splatting and Monocular Reconstruction}

3D Gaussian Splatting (3DGS)~\cite{kerbl20233dgs} represents a scene as a
collection of anisotropic 3D Gaussians with differentiable rasterization,
enabling real-time photorealistic rendering. 2D Gaussian Splatting
(2DGS)~\cite{huang20242dgs} collapses the third scale dimension into a surface
normal, improving geometric fidelity for surface-like content.
Feed-forward monocular reconstruction models have recently emerged that
predict 3DGS parameters directly from a single image:
DreamGaussian~\cite{tang2024dreamgaussian} uses a generative prior for
efficient 3D content creation; TripoSplat~\cite{wang2024triposplat} achieves
state-of-the-art single-image-to-3DGS quality; pixelSplat~\cite{charatan2024pixelsplat}
and AnySplat~\cite{jiang2025anysplat} handle two-view and unconstrained-view
settings respectively. These advances make it feasible to lift individual
painting regions into 3D---but none address the downstream challenge of
composing multiple such objects into a coherent scene with artistic control.
Canvas3D leverages TripoSplat per-object and addresses the composition and
manipulation gap.

\subsection{Hand-Drawn Sketch and Line-Art Guidance}

Hand-drawn sketches have long served as an intuitive interface for visual
content creation. In 3D modeling, sketch-based interfaces~\cite{olson2011sketch}
enable rapid shape prototyping. ShadowDraw~\cite{luo2026shadowdraw} generates
compositional shadow-drawing art from object inputs, demonstrating that
contour-based spatial reasoning is tractable with modern vision models.
LineArt~\cite{wang2024lineart} performs training-free appearance transfer onto
design drawings via diffusion models. For 3D scene composition, region-aware
line drawing simplification~\cite{nguyen2024region} and volumetric animation
toolkits~\cite{foxgieg2024lightning} illustrate the expressive potential of
hand-drawn control. Canvas3D extends this lineage by using line-art not merely
for stylization, but as the \textit{spatial constraint} that solves the
which-object-where problem in 3D scene assembly.

\subsection{Controlled Image-to-Video Generation}

The leap from text-to-image to text-to-video has been powered by diffusion
transformers~(DiT)~\cite{peebles2023dit} and video diffusion
models~\cite{ho2022vdm}. CogVideoX~\cite{yang2025cogvideox} introduces an expert
transformer for efficient video generation. Wan2.2~\cite{wan2025} scales to
large MoE architectures with text- and image-conditioned video synthesis.
LTX-Video~\cite{ltxvideo2024} achieves real-time generation with aggressive
VAE compression.

Controllability has been addressed through auxiliary conditioning.
ControlNet~\cite{zhang2023controlnet} injects edge, depth, and pose maps into
the diffusion U-Net; EasyControl~\cite{wang2024easycontrol} extends this to
U-Net--based video diffusion; ControlNeXt~\cite{peng2024controlnext} reduces
parameter overhead by up to 90\%. MotionCanvas~\cite{xing2025motioncanvas}
achieves cinematic shot design with 3D-aware motion control in I2V without
3D training data. AniClipart~\cite{wu2025aniclipart} animates clipart via
Video Score Distillation Sampling from text-to-video priors.

Despite these advances, existing I2V control methods are limited to
\textit{per-frame conditioning signals}---they cannot reason about 3D object
identity, depth ordering, or occlusion across frames. Canvas3D addresses
this by grounding all control signals in an explicit 3D scene representation
that is composited, rendered, and then fed into the video diffusion pipeline.

\subsection{Artistic Video Generation and Painting Animation}

Dedicated painting animation methods span several paradigms.
Stroke-level synthesis~\cite{zou2021stylized,liu2021neuralstrokes}
reproduces the brushstroke generation process as a time-lapse but does not
introduce novel motion. Portrait animation methods such as
LivePortrait~\cite{guo2024liveportrait} and EDTalk~\cite{tan2024edtalk}
drive facial motion on painted faces but cannot animate full-body figures
or background elements. The recently proposed ``Printing the
Underdetermined''~\cite{ming2026printing} translates paintings to 3DGS via video
generative models but operates on the entire image jointly rather than at the
object level. OilPaintingRevive~\cite{oilpaintingrevive2026} proposes a
five-stage pipeline combining 2DGS with LTX-Video refinement, but its skeleton
binding network requires per-scene optimization that limits compositional
flexibility.

Canvas3D advances beyond these works by: (i)~operating at the object level
rather than the image level, enabling independent manipulation; (ii)~using
hand-drawn line-art as the compositional control interface, separating artistic
intent from algorithmic prediction; (iii)~integrating depth-aware inpainting to
handle disocclusion; and (iv)~employing multi-modal ControlNet (depth~+~edge)
for geometry-grounded video synthesis with stronger 3D consistency guarantees.
